\documentclass{article}
\usepackage[utf8]{inputenc}
\usepackage[brazil]{babel}
\usepackage{geometry}
\usepackage{titlesec}
\usepackage{enumitem}
\usepackage{hyperref}
\usepackage{xcolor}
\usepackage{setspace}
\usepackage{fancyhdr}
\usepackage{listings}

\geometry{margin=2.5cm}

\titleformat{\section}{\Large\bfseries}{\thesection}{1em}{}
\titleformat{\subsection}{\bfseries}{\thesubsection}{1em}{}

\lstset{
    basicstyle=\ttfamily\small,
    backgroundcolor=\color{gray!10},
    frame=single,
    breaklines=true,
    showstringspaces=false
}

\begin{document}

\begin{center}
    \centering

    {\huge \textbf{Análise, Projeto e Desenvolvimento Ágil}}\\[0.5cm]

    {\Huge \textbf{Documento Instrucional para RH}}\\[0.5cm]

    {\large Aluno: \textbf{João Miguel de Castro Menna}}\\[0.5cm]
\end{center}

\vspace{1em}
\hrule
\vspace{1em}

\section{Resumo Executivo}

\paragraph{}

Este documento foi elaborado para auxiliar o setor de RH na triagem, entrevistas e onboarding de novos(as) desenvolvedores(as), mesmo sem conhecimento prévio em metodologias ágeis.

A empresa adota simultaneamente três abordagens: \textbf{Scrum}, \textbf{Kanban} e \textbf{Extreme Programming (XP)}. Cada uma é aplicada de forma integrada, garantindo previsibilidade, fluxo contínuo e qualidade técnica sustentada por testes e refatoração.

O documento também descreve o perfil esperado, o processo de seleção e exemplos práticos que podem ser utilizados em entrevistas ou treinamentos.

\section{Descrição da Vaga e Pré-Requisitos}

\textbf{Título:} Desenvolvedor(a) de Software Ágil - Pleno/Sênior

\subsection*{Pré-requisitos obrigatórios}
\begin{itemize}[noitemsep]
    \item Experiência prática com \textbf{Scrum, Kanban e XP/TDD}.
    \item Participação em equipes ágeis, entregando incrementos de software de forma iterativa e contínua.
\end{itemize}

\subsection*{Hard Skills}
\begin{itemize}[noitemsep]
    \item Testes automatizados (unitários, integração, end-to-end);
    \item Versionamento Git (fluxos com \textit{feature branches});
    \item Integração Contínua (CI) e entrega contínua (CD);
    \item Refatoração e revisão de código.
\end{itemize}

\subsection*{Soft Skills}
\begin{itemize}[noitemsep]
    \item Colaboração e trabalho em equipe;
    \item Comunicação clara e empática;
    \item Feedbacks curtos e construtivos;
    \item Foco em qualidade e melhoria contínua.
\end{itemize}

\section{Como Trabalhamos com Scrum}

\subsection*{Papéis}
\begin{itemize}[noitemsep]
    \item \textbf{Product Owner (PO):} Prioriza o backlog e garante alinhamento com o valor de negócio.
    \item \textbf{Scrum Master (SM):} Remove impedimentos e assegura a aplicação correta do processo.
    \item \textbf{Dev Team:} Entrega um incremento potencialmente liberável a cada Sprint.
\end{itemize}

\subsection*{Eventos (Time-box de referência)}
\begin{itemize}[noitemsep]
    \item \textbf{Sprint Planning (até 4h / 2 semanas):} Definição da meta da Sprint e seleção dos itens priorizados.
    \item \textbf{Daily (15 min):} Sincronização do time, verificação do quadro Kanban e identificação de impedimentos.
    \item \textbf{Review (até 2h):} Apresentação do incremento aos stakeholders.
    \item \textbf{Retrospective (até 1,5h):} Reflexão sobre o processo e definição de 1–2 ações de melhoria.
\end{itemize}

\subsection*{Artefatos e Definições}
\begin{itemize}[noitemsep]
    \item \textbf{Product Backlog} e \textbf{Sprint Backlog} com \textit{Definition of Ready} (DoR) e \textit{Definition of Done} (DoD);
    \item \textbf{Incremento} testado, integrado e potencialmente liberável.
\end{itemize}

\section{Nosso Kanban Integrado ao Scrum}

O Kanban é utilizado para visualizar e gerenciar o fluxo de trabalho dentro da Sprint.

\subsection*{Exemplo de Quadro Kanban}
\begin{center}
Backlog → Ready → In Progress (WIP 3) → Code Review (WIP 2) → Test (WIP 2) → Done
\end{center}

\subsection*{Políticas Explícitas}
\begin{itemize}[noitemsep]
    \item Cada coluna possui critérios de entrada e saída claros;
    \item Itens bloqueados são sinalizados com etiqueta “Blocked” e motivo;
    \item Pull só é permitido quando a coluna atual está abaixo do limite de WIP;
    \item Itens urgentes seguem a classe especial “Expedite” (WIP dedicado).
\end{itemize}

\subsection*{Métricas}
\begin{itemize}[noitemsep]
    \item \textbf{Lead Time:} tempo total do item até a entrega;
    \item \textbf{Throughput:} quantidade média de entregas semanais.
\end{itemize}

\section{XP e TDD na Prática}

O time adota práticas de \textbf{Extreme Programming (XP)} para garantir qualidade e sustentabilidade do código.

\subsection*{Ciclo TDD (Red–Green–Refactor)}
\begin{enumerate}
    \item \textbf{Red:} escrever um teste que falha.
    \item \textbf{Green:} implementar o código mínimo para o teste passar.
    \item \textbf{Refactor:} melhorar o código mantendo todos os testes verdes.
\end{enumerate}

\subsection*{Práticas Complementares}
\begin{itemize}[noitemsep]
    \item \textbf{Pair Programming} (ou mob em casos críticos);
    \item \textbf{Code Review} leve e colaborativo;
    \item \textbf{Integração Contínua (CI)} a cada \textit{push};
    \item \textbf{Feature Flags} para entregas graduais;
    \item \textbf{Commits pequenos} e \textbf{refatoração contínua}.
\end{itemize}

\section{Roteiro de Seleção para RH}

\subsection*{Triagem de Currículo}
Verificar menções a:
\begin{itemize}[noitemsep]
    \item TDD, testes automatizados e refatoração;
    \item Experiência com Scrum e/ou Kanban;
    \item Práticas de CI/CD.
\end{itemize}

\subsection*{Perguntas para Entrevista (linguagem simples)}
\begin{itemize}[noitemsep]
    \item “Explique, em linguagem simples, o que é TDD e por que usamos.”
    \item “O que significa limite de WIP e por que ele ajuda?”
    \item “Como você participa de uma Daily produtiva?”
\end{itemize}

\subsection*{Critérios de Avaliação Técnica}
\begin{itemize}[noitemsep]
    \item Mini-desafio técnico com TDD;
    \item Leitura e explicação de um trecho de código real;
    \item Simulação de uma revisão de código (code review);
    \item Discussão de caso sobre priorização de backlog.
\end{itemize}

\section{Exemplos Práticos}

\subsection*{User Story (Scrum)}
\textbf{Como} usuário autenticado,  
\textbf{quero} redefinir minha senha via link seguro,  
\textbf{para} poder recuperar o acesso rapidamente.  

\textbf{Critérios de Aceite:}
\begin{itemize}[noitemsep]
    \item O link de redefinição expira em 30 minutos;
    \item O usuário recebe um e-mail com o link seguro;
    \item Após redefinir, o token é invalidado.
\end{itemize}

\subsection*{Quadro Kanban (resumo visual)}
\begin{center}
\texttt{Backlog → Ready → In Progress (WIP 3) → Code Review (WIP 2) → Test (WIP 2) → Done}
\end{center}

\subsection*{Mini Ciclo TDD (Pseudo-código)}
\begin{lstlisting}[language=Python]
# Red
assert(reset_link.expires_in_minutes == 30)  # deve falhar inicialmente

# Green
reset_link.set_expiration(30)  # torna o teste verde

# Refactor
DEFAULT_RESET_EXPIRATION = 30
# extrair funcao, remover duplicacoes, manter testes passando
\end{lstlisting}

\section{Anexos e Diagramas (opcional)}
\begin{itemize}[noitemsep]
    \item Fluxo resumido: Backlog → Sprint → Kanban → Entrega → Feedback.
    \item Diagrama simples de papéis (PO, SM, Dev Team).
    \item Exemplo de métricas (gráfico de Lead Time e Throughput).
\end{itemize}

\vfill
\begin{center}
\textit{“A cultura ágil é construída todos os dias — começando pela forma como contratamos e integramos pessoas.”}
\end{center}

\end{document}
